\section{Conclusion}

This study establishes that ordinal crop rankings do not, by themselves, determine constrained land allocations. A set-valued stochastic programme clarifies possible, universal and selected reversal, while repaired structural results show how feasibility, downside risk, dependence and optimizer multiplicity enter without implying a universal threshold. Official state data document substantial but definition-sensitive ranking--acreage discordance, alongside a national null. The simulation provides transparent diagnostics but remains nonheadline because its convergence rule fails.

The combined evidence supports a bounded conclusion: ranking-to-acreage disagreement is real in the selected descriptive sample, whereas the optimizing mechanism behind observed acreage is not identified. Longer panels, farm-level constraints and expectations, and adequately replicated preregistered simulations are required to distinguish operational, risk and dependence mechanisms. Until then, rankings are best treated as inputs to explicit decisions rather than as allocations in disguise.
